% GENERATED FILE. Edit canonical lessons, then run npm run generate:books.
% canonical-source-hash: fnv1a64:a30b176eb40c0a19
% canonical-lessons: MR-C489-las, MR-C489-sastrakriya, MR-C489-bhul, MR-C489-taka, MR-C489-tocne, MR-R489-review-482-489

\chapter{Needles and Stitches}
\label{ch:mr-needles-and-stitches}
\hlchaptermodality{489}

\begin{chapteropening}
\textbf{By the end of this chapter:} \emph{I can talk about a vaccine, an operation, anaesthesia and stitches.}

Complete the last lesson of chapter 489: I can talk about a vaccine, an operation, anaesthesia and stitches.
\end{chapteropening}

\section[\texorpdfstring{\mr{लस} (las)}{लस (las)}]{\mr{लस} (las) — a vaccine}

\label{lesson:MR-C489-las}



\begin{quote}
Before the new one: say the Marathi for a diagnosis, then the Marathi for treatment.
\end{quote}

\subsection*{You'll want to know: \mr{लस}}
\textbf{\mr{लस}} — \emph{las} — \textquotedblleft{}a vaccine\textquotedblright{}.

\begin{quote}
\mr{लहान} \mr{बाळांना} \mr{वेळेवर} \mr{लस} \mr{द्यायला} \mr{हवी}.
\end{quote}

\begin{quote}
\emph{lahān bāḷānnā veḷevar las dyāylā havī.}
\end{quote}

\begin{quote}
Small babies should be vaccinated on time.
\end{quote}

Feminine.

\mr{लस} \mr{देणे}, to vaccinate; \mr{लस} \mr{घेणे}, to get vaccinated.

\begin{cousinweb}
\textbf{\mr{लस}} (\emph{las}) is inherited: CDIAL lists Marathi \emph{las}, \textquotedblleft{}serum, pus\textquotedblright{}, under the reconstructed \emph{*lasa}, \textquotedblleft{}stickiness\textquotedblright{} (CDIAL 10992). The sense \textquotedblleft{}vaccine\textquotedblright{} is modern, usually explained by the serum vaccines were made from.
\end{cousinweb}

\begin{grammarlens}[title={What you've built}]
One more word for describing what you lived and felt.
\end{grammarlens}

\subsection*{Your turn}
\begin{itemize}
\raggedright
  \item \emph{Say it:} \emph{las}
  \item \emph{Say it:} \emph{las}, once more
  \item \emph{Say it:} \emph{lahān bāḷānnā veḷevar las dyāylā havī.} — Small babies should be vaccinated on time.
  \item \emph{Recall:} say the Marathi for a diagnosis, then the Marathi for treatment, then say \emph{las} again
\end{itemize}

\begin{tcolorbox}[breakable,colback=teal!4,colframe=teal!35!black,title={Before you move on}]
What does \mr{लस} mean? (\textquotedblleft{}A vaccine\textquotedblright{}.) Say it once more.
\end{tcolorbox}
\section[\texorpdfstring{\mr{शस्त्रक्रिया} (śastrakriyā)}{शस्त्रक्रिया (śastrakriyā)}]{\mr{शस्त्रक्रिया} (śastrakriyā) — an operation (surgery)}

\label{lesson:MR-C489-sastrakriya}



\begin{quote}
Before the new one: say the Marathi for treatment, then the Marathi for a vaccine.
\end{quote}

\subsection*{You'll want to know: \mr{शस्त्रक्रिया}}
\textbf{\mr{शस्त्रक्रिया}} — \emph{śastrakriyā} — \textquotedblleft{}an operation (surgery)\textquotedblright{}.

\begin{quote}
\mr{आजोबांच्या} \mr{डोळ्याची} \mr{शस्त्रक्रिया} \mr{यशस्वी} \mr{झाली}.
\end{quote}

\begin{quote}
\emph{ājobāñcyā ḍoḷyācī śastrakriyā yaśasvī jhālī.}
\end{quote}

\begin{quote}
Grandfather's eye operation was a success.
\end{quote}

Feminine.

\mr{शस्त्र}, a blade, + \mr{क्रिया}, an act.

\begin{cousinweb}
\textbf{\mr{शस्त्रक्रिया}} (\emph{śastrakriyā}) is a modern learned compound built from two Sanskrit words, \emph{śastra}, \textquotedblleft{}knife, blade\textquotedblright{}, and \emph{kriyā}, \textquotedblleft{}doing, act\textquotedblright{} (Monier-Williams).
\end{cousinweb}

\begin{grammarlens}[title={What you've built}]
One more word for describing what you lived and felt.
\end{grammarlens}

\subsection*{Your turn}
\begin{itemize}
\raggedright
  \item \emph{Say it:} \emph{śastrakriyā}
  \item \emph{Say it:} \emph{śastrakriyā}, once more
  \item \emph{Say it:} \emph{ājobāñcyā ḍoḷyācī śastrakriyā yaśasvī jhālī.} — Grandfather's eye operation was a success.
  \item \emph{Recall:} say the Marathi for treatment, then the Marathi for a vaccine, then say \emph{śastrakriyā} again
\end{itemize}

\begin{tcolorbox}[breakable,colback=teal!4,colframe=teal!35!black,title={Before you move on}]
What does \mr{शस्त्रक्रिया} mean? (\textquotedblleft{}An operation (surgery)\textquotedblright{}.) Say it once more.
\end{tcolorbox}
\section[\texorpdfstring{\mr{भूल} (bhūl)}{भूल (bhūl)}]{\mr{भूल} (bhūl) — anaesthesia}

\label{lesson:MR-C489-bhul}



\begin{quote}
Before the new one: say the Marathi for a vaccine, then the Marathi for an operation.
\end{quote}

\subsection*{You'll want to know: \mr{भूल}}
\textbf{\mr{भूल}} — \emph{bhūl} — \textquotedblleft{}anaesthesia\textquotedblright{}.

\begin{quote}
\mr{शस्त्रक्रियेआधी} \mr{त्याला} \mr{भूल} \mr{दिली}.
\end{quote}

\begin{quote}
\emph{śastrakriyeādhī tyālā bhūl dilī.}
\end{quote}

\begin{quote}
He was given anaesthesia before the operation.
\end{quote}

Feminine.

\mr{भूल} \mr{देणे}, to anaesthetise. The same word means a slip of memory.

\begin{cousinweb}
\textbf{\mr{भूल}} (\emph{bhūl}) is inherited: CDIAL lists Marathi \emph{bhulṇẽ}, \textquotedblleft{}errs, forgets\textquotedblright{}, under the reconstructed \emph{*bhull} (CDIAL 9538); \mr{भूल} is the noun, and modern Marathi gave it the sense \textquotedblleft{}anaesthesia\textquotedblright{}.
\end{cousinweb}

\begin{grammarlens}[title={What you've built}]
One more word for describing what you lived and felt.
\end{grammarlens}

\subsection*{Your turn}
\begin{itemize}
\raggedright
  \item \emph{Say it:} \emph{bhūl}
  \item \emph{Say it:} \emph{bhūl}, once more
  \item \emph{Say it:} \emph{śastrakriyeādhī tyālā bhūl dilī.} — He was given anaesthesia before the operation.
  \item \emph{Recall:} say the Marathi for a vaccine, then the Marathi for an operation, then say \emph{bhūl} again
\end{itemize}

\begin{tcolorbox}[breakable,colback=teal!4,colframe=teal!35!black,title={Before you move on}]
What does \mr{भूल} mean? (\textquotedblleft{}Anaesthesia\textquotedblright{}.) Say it once more.
\end{tcolorbox}
\section[\texorpdfstring{\mr{टाका} (ṭākā)}{टाका (ṭākā)}]{\mr{टाका} (ṭākā) — a stitch (in a wound)}

\label{lesson:MR-C489-taka}



\begin{quote}
Before the new one: say the Marathi for an operation, then the Marathi for anaesthesia.
\end{quote}

\subsection*{You'll want to know: \mr{टाका}}
\textbf{\mr{टाका}} — \emph{ṭākā} — \textquotedblleft{}a stitch (in a wound)\textquotedblright{}.

\begin{quote}
\mr{जखमेला} \mr{चार} \mr{टाके} \mr{घातले}.
\end{quote}

\begin{quote}
\emph{jakhmelā cār ṭāke ghātle.}
\end{quote}

\begin{quote}
They put four stitches in the wound.
\end{quote}

Masculine. Plural: \mr{टाके}, \emph{ṭāke}.

Usually plural, \mr{टाके}: \mr{टाके} \mr{घालणे}, to put in stitches; \mr{टाके} \mr{काढणे}, to take them out.

\begin{cousinweb}
\textbf{\mr{टाका}} (\emph{ṭākā}) is inherited: CDIAL lists Marathi \emph{ṭākā}, of stitching, under Sanskrit \emph{ṭaṅkati}, \textquotedblleft{}ties\textquotedblright{} (CDIAL 5432).
\end{cousinweb}

\begin{grammarlens}[title={What you've built}]
One more word for describing what you lived and felt.
\end{grammarlens}

\subsection*{Your turn}
\begin{itemize}
\raggedright
  \item \emph{Say it:} \emph{ṭākā}
  \item \emph{Say it:} \emph{ṭākā}, once more
  \item \emph{Say it:} \emph{jakhmelā cār ṭāke ghātle.} — They put four stitches in the wound.
  \item \emph{Recall:} say the Marathi for an operation, then the Marathi for anaesthesia, then say \emph{ṭākā} again
\end{itemize}

\begin{tcolorbox}[breakable,colback=teal!4,colframe=teal!35!black,title={Before you move on}]
What does \mr{टाका} mean? (\textquotedblleft{}A stitch (in a wound)\textquotedblright{}.) Say it once more.
\end{tcolorbox}
\section[\texorpdfstring{\mr{टोचणे} (ṭocṇe)}{टोचणे (ṭocṇe)}]{\mr{टोचणे} (ṭocṇe) — to prick; to inject}

\label{lesson:MR-C489-tocne}



\begin{quote}
Before the new one: say the Marathi for anaesthesia, then the Marathi for a stitch.
\end{quote}

\subsection*{You'll want to know: \mr{टोचणे}}
\textbf{\mr{टोचणे}} — \emph{ṭocṇe} — \textquotedblleft{}to prick; to inject\textquotedblright{}.

\begin{quote}
\mr{परिचारिकेने} \mr{माझ्या} \mr{हातावर} \mr{सुई} \mr{टोचली}.
\end{quote}

\begin{quote}
\emph{paricārikene mājhyā hātāvar suī ṭoclī.}
\end{quote}

\begin{quote}
The nurse pricked my arm with a needle.
\end{quote}

Past, with -\mr{ने} (he …): \mr{त्याने} \mr{सुई} \mr{टोचली}, \emph{tyāne suī ṭoclī}.

Of thorns, needles and injections: \mr{काटा} \mr{टोचला}, a thorn pricked.

\begin{cousinweb}
\textbf{\mr{टोचणे}} (\emph{ṭocṇe}) is a Marathi verb whose origin is not settled; CDIAL lists a different Marathi verb of pricking, \emph{ṭoviṇẽ}, under \emph{trōpati} (CDIAL 6083).
\end{cousinweb}

\begin{grammarlens}[title={What you've built}]
One more word for describing what you lived and felt.
\end{grammarlens}

\subsection*{Your turn}
\begin{itemize}
\raggedright
  \item \emph{Say it:} \emph{ṭocṇe}
  \item \emph{Say it:} \emph{ṭocṇe}, once more
  \item \emph{Say it:} \emph{paricārikene mājhyā hātāvar suī ṭoclī.} — The nurse pricked my arm with a needle.
  \item \emph{Recall:} say the Marathi for anaesthesia, then the Marathi for a stitch, then say \emph{ṭocṇe} again
\end{itemize}

\begin{tcolorbox}[breakable,colback=teal!4,colframe=teal!35!black,title={Before you move on}]
What does \mr{टोचणे} mean? (\textquotedblleft{}To prick; to inject\textquotedblright{}.) Say it once more.
\end{tcolorbox}
\section[(dialogue)]{Looking back: from Aches and Pains to Needles and Stitches}

\label{lesson:MR-R489-review-482-489}



\begin{quote}
Say the Marathi for a headache and to ache.
\end{quote}

\subsection*{Your turn}
Say these aloud:

\begin{itemize}
\raggedright
  \item a headache, to ache, to get a twinge, to groan, to moan in pain
  \item indigestion, acidity, diarrhoea, to vomit, to be digested
  \item a disease, an infection, a germ, a virus, an epidemic
  \item diabetes, blood pressure, asthma, jaundice, chickenpox
  \item an injury, a scab, to get grazed, to sprain, to swell
  \item to get a burn, an abscess, a rash, to roast, to scratch
  \item a sick person, a nurse, a physician, a diagnosis, treatment
  \item a vaccine, an operation, anaesthesia, a stitch, to prick
\end{itemize}

\begin{tcolorbox}[breakable,colback=teal!4,colframe=teal!35!black,title={Before you move on}]
A headache? (\textbf{\mr{डोकेदुखी}}, \emph{ḍokedukhī}.) To ache? (\textbf{\mr{दुखणे}}, \emph{dukhṇe}.) To get a twinge? (\textbf{\mr{कळ} \mr{येणे}}, \emph{kaḷ yeṇe}.) To groan? (\textbf{\mr{कण्हणे}}, \emph{kaṇhṇe}.) To moan in pain? (\textbf{\mr{विव्हळणे}}, \emph{vivhaḷṇe}.) Indigestion? (\textbf{\mr{अपचन}}, \emph{apcan}.) Acidity? (\textbf{\mr{पित्त}}, \emph{pitta}.) Diarrhoea? (\textbf{\mr{जुलाब}}, \emph{julāb}.) To vomit? (\textbf{\mr{ओकणे}}, \emph{okṇe}.) To be digested? (\textbf{\mr{पचणे}}, \emph{pacṇe}.) A disease? (\textbf{\mr{रोग}}, \emph{rog}.) An infection? (\textbf{\mr{संसर्ग}}, \emph{saṁsarga}.) A germ? (\textbf{\mr{जंतू}}, \emph{jantū}.) A virus? (\textbf{\mr{विषाणू}}, \emph{viṣāṇū}.) An epidemic? (\textbf{\mr{साथ}}, \emph{sāth}.) Diabetes? (\textbf{\mr{मधुमेह}}, \emph{madhumeh}.) Blood pressure? (\textbf{\mr{रक्तदाब}}, \emph{raktadāb}.) Asthma? (\textbf{\mr{दमा}}, \emph{damā}.) Jaundice? (\textbf{\mr{कावीळ}}, \emph{kāvīḷ}.) Chickenpox? (\textbf{\mr{कांजिण्या}}, \emph{kāñjiṇyā}.) An injury? (\textbf{\mr{इजा}}, \emph{ijā}.) A scab? (\textbf{\mr{खपली}}, \emph{khaplī}.) To get grazed? (\textbf{\mr{खरचटणे}}, \emph{kharcaṭṇe}.) To sprain? (\textbf{\mr{मुरगळणे}}, \emph{murgaḷṇe}.) To swell? (\textbf{\mr{सुजणे}}, \emph{sujṇe}.) To get a burn? (\textbf{\mr{चटका} \mr{बसणे}}, \emph{caṭkā basṇe}.) An abscess? (\textbf{\mr{गळू}}, \emph{gaḷū}.) A rash? (\textbf{\mr{पुरळ}}, \emph{puraḷ}.) To roast? (\textbf{\mr{भाजणे}}, \emph{bhājṇe}.) To scratch? (\textbf{\mr{खाजवणे}}, \emph{khājavṇe}.) A sick person? (\textbf{\mr{रुग्ण}}, \emph{rugṇa}.) A nurse? (\textbf{\mr{परिचारिका}}, \emph{paricārikā}.) A physician? (\textbf{\mr{वैद्य}}, \emph{vaidya}.) A diagnosis? (\textbf{\mr{निदान}}, \emph{nidān}.) Treatment? (\textbf{\mr{उपचार}}, \emph{upcār}.) A vaccine? (\textbf{\mr{लस}}, \emph{las}.) An operation? (\textbf{\mr{शस्त्रक्रिया}}, \emph{śastrakriyā}.) Anaesthesia? (\textbf{\mr{भूल}}, \emph{bhūl}.) A stitch? (\textbf{\mr{टाका}}, \emph{ṭākā}.) To prick? (\textbf{\mr{टोचणे}}, \emph{ṭocṇe}.)
\end{tcolorbox}
